\section{Chapter 4}

\begin{quotex}
\begin{verse}
The Tao is not substance, it is inexhaustible activity\\ Acting does not increase it\\ As it is unfathomable!\\ It is like the primordial source of all things.\\ It blunts the sharp\\ It clarifies the confused\\ It dims the bright\\ It orders the elementary parts ``powder'' (of matter)\\ It is elusive and yet omnipresent\\ How could it have been begotten?\\ It is anterior (and superior) to the Lord of Heaven
\end{verse}
\end{quotex}


\paragraph{Commentary}


Insubstantiality (``emptiness'') is emphasized, the pure activity of the Principle, that neither expands nor diminishes it (therefore in opposition both to the theory of ``becoming'' of the Spirit, and to that of its deterioration through emanation). The Taoist image is also that of a vase that, while turning continually, what it contains remains full and in order to fill itself, still remains empty. ``Primordial source'' is the abstract meaning of \emph{tsung} in the text, which is literally ancestor, the original father. If the image is conserved, given the part of the regulator typical of the Chinese head of the family, in the fourth line, following some translations, the idea of an ordering power can also be included. Then a natural connection is established with the four successive lines that are about the action of the Tao.

For Taoist metaphysic, the two last lines are important, where the principle is called ``without origin'' --- \emph{wu yuan} --- superior and anterior to the King of Heaven (the Shangdi\footnote{\url{https://en.wikipedia.org/wiki/Shangdi}}), i.e., to the personal God. In this passage a certain circumspect way of expressing itself was emphasized linguistically, due to the part that the Shangdi had in the cult of the State under Chou's dynasty, when the Tao Te Ching was composed. It is like that through a concern for that conception of exoterism. But in the Introduction, we noted the irrelevance of the personal, theistic and anthropomorphic conceptions of the divine to the primordial tradition of the Far East, taken up in its abstract, metaphysical orientation by Lao Tzu. Chuang Tzu (II,3) makes clear that one can admit the principle of a universal regulator (from which the principle of every unity, family, people, etc. derives through participation) but on condition that one does not make a distinct personal being out of what is intended ``as an influence without a comprehensible form''.

\paragraph{Chinese text and literal translation}


Chapter 4 (\begin{CJK*}{UTF8}{bsmi}第四章\end{CJK*})
\columnratio{0.4}
\begin{paracol}{2}
\begin{verse}
\begin{CJK*}{UTF8}{bsmi}
道沖，\\
而用之或不盈。\\ 淵兮，似萬物之宗。\\ 挫其銳，\\ 解其紛，\\ 和其光，\\ 同其塵。\\ 湛兮，似或存。\\ 吾不知誰之子，\\ 象帝之先。
\end{CJK*}
\end{verse}
\switchcolumn
\begin{verse}
The Dao is empty,\\ But when using it, it is impossible to use it up.\\ It is profound, seems like the root of the myriad things.\\ Blunts its own sharpness.\\ Unravels its own fetters.\\ Harmonises its own light.\\ Mixes with its own dust.\\ It is unclear, but seems to have existed there.\\ I do not know whose son it is,\\ Maybe it was already created before the creator.
\end{verse}
\end{paracol}

\flrightit{Posted on 2020-04-09 by Lao Tzu}

